
\section{Equivalent-electron terms: \(nd^2\)}
Each \(d\) electron has \(l=2\), \(s=\frac12\), so ten
spin-orbitals are available. Two equivalent electrons
have \(\binom{10}{2}=45\) antisymmetric microstates.
The allowed LS terms are
\[
{}^3F,\quad{}^3P,\quad{}^1G,\quad{}^1D,\quad{}^1S.
\]
Split into \(J\) values:
\[
{}^3F_{2,3,4},\quad{}^3P_{0,1,2},\quad
{}^1G_4,\quad{}^1D_2,\quad{}^1S_0.
\]
The check is \(21+9+9+5+1=45\).
Hund's rules suggest a \({}^3F_2\) ground level for less-than-half-filled
\(d^2\) in the LS-coupled limit.
\begin{methodbox}
\textbf{Method:} count microstates \(\binom{4l+2}{n}\);
couple \(L,S\), enforce exchange antisymmetry, list \(J\),
sum \((2J+1)\), then state which Hund-rule comparisons are justified.
Compare observed levels for the chosen element without claiming
the rules determine every excited-term energy.
\end{methodbox}
